% answers_short.tex —— 1977 年·江苏盐城地区初试卷 · 简短答案速查
% 由 tools/make_short_answers.py 自 answers.tex 的 \gksolhead/\gkanswer 提取，
% 并与 1977/1977_数学真题_江苏盐城地区初试_答案要点.md 逐题比对后生成，勿手改（改源 MD 后重跑）。

\begin{gknotice}
\textbf{使用说明}：本节只给最终答案便于快速核对，完整推导见解析版。答案为编写者独立推导，\textbf{非官方标准答案}。
\end{gknotice}

\gkpart{题目}
\begin{itemize}[leftmargin=2.6em]
  \item[\textbf{1.}] (1) $-\dfrac13$；(2) $-\dfrac12$
  \item[\textbf{2.}] (1) $-5\le x\le12$；(2) $y=5\sqrt2$
  \item[\textbf{3.}] $m=1$ 时，方程有两个相等的实数根 $x_1=x_2=1$。
  \item[\textbf{4.}] $\odot O$ 的面积为 $\dfrac{80\pi}{13}\ \mathrm{cm^2}$。
  \item[\textbf{5.}] $\dfrac{x^2}5+\dfrac{y^2}3=1$。
  \item[\textbf{6.}] 见解析版
  \item[\textbf{7.}] (1) 当 $A_1\text{、}B_1\text{、}C_1\text{、}D_1$ 分别是正方形 $ABCD$ 各边中点时，内接正方形 $A_1B_1C_1D_1$ 的面积最小，最小值为 $\dfrac12$；(2) 第 $n$ 个面积最小的内接正方形的面积为 $\dfrac1{2^n}$；(3) $1$
  \item[\textbf{8.}] 见解析版
\end{itemize}
